% Teil II: der durchgerechnete Weg.

% Gedankenschritt: die Formeln aus dem Flächenbild und aus dem
% Distributivgesetz verstehen, statt sie nur nachzuschlagen.

\subsection*{Gegeben und gesucht}
Ausmultiplizieren von $(x+y)^2$, $(x-y)^2$ und $(x+y)(x-y)$ für Variablen
$x$, $y$.

\subsection*{Ansatz und der entscheidende Gedankenschritt}
Statt die Formeln nur nachzuschlagen, rekonstruieren wir sie aus dem
Distributivgesetz --- das ist der Schlüssel, um sie wirklich ``zu haben``.

\textbf{Vorgehen:} $(x+y)^2$ schreiben wir als $(x+y)(x+y)$ und wenden das
Distributivgesetz an: jeder von beiden Summanden im ersten Faktor
multipliziert beide Summanden im zweiten. Das liefert vier Produkte;
von diesen fassen wir die beiden ``gemischten`` ($xy$ und $yx$) zusammen.

\subsection*{Rechnung}
\begin{align}
  (x+y)^2 &= x^2 + xy + yx + y^2 = x^2+2xy+y^2, \\
  (x-y)^2 &= (x-y)(x-y)=x^2-xy-yx+y^2 = x^2-2xy+y^2, \\
  (x+y)(x-y) &= x^2-xy+yx-y^2 = x^2-y^2 .
\end{align}
Bei den beiden ersten zählen wir $xy+yx=2$ gemischte Produkte (eines mit
$+$, eines mit $-$, je nach Vorzeichen). Bei der dritten heben sich die
beiden Mittelterme $xy$ und $yx$ mit \emph{entgegengesetztem} Vorzeichen
weg.

\begin{ergebnis}
\begin{align*}
  (x+y)^2 &= x^2+2xy+y^2, \\
  (x-y)^2 &= x^2-2xy+y^2, \\
  (x+y)(x-y) &= x^2-y^2 .
\end{align*}
\end{ergebnis}

\begin{kontrolle}
\textbf{Probe mit Zahlen (eine starke Kontrolle):} setzt man $x=1,
y=1$: links $(1+1)^2=4$, rechts $1+2+1=4$; $(1-1)^2=0$, rechts $1-2+1=0$;
$(1+1)(1-1)=0$, rechts $1-1=0$. In allen drei Fällen stimmt die Probe.

\textbf{Grenzfall $y=0$:} Setzt man $y=0$ ein, wird jede der drei Formeln
zur trivialen Identität $x^2=x^2$ --- ein guter ``Selbsttest`` der
Struktur, der zeigt, dass das doppelte Produkt logisch zusammengehört.
Hätte der Mittelterm $2xy$ nicht gestimmt, wäre diese Probe für $y=0$
zwar harmlos (weil $2xy=0$), aber für $y=1$ sofort sichtbar: Es bleibt
$x^2+1$ statt $(x+1)^2=x^2+2x+1$, wenn man $2x$ vergisst.
\end{kontrolle}

\begin{merksatz}
Die binomischen Formeln sind Spezialfälle des Distributivgesetzes; der
Mittelterm $2xy$ (bzw. $-2xy$) ist beim Quadrieren nicht zu vergessen ---
anschaulich entstehen vier Teilflächen, von denen die beiden Rechtecke
das doppelte Produkt liefern.
\end{merksatz}
